% =====================================================================
%  1bat-ud1-6.tex  —  HORA 6: El parany del signe negatiu (resoldre inequacions)
% =====================================================================
\horatitol{Sessió 6 — El parany del signe negatiu: com es resolen les inequacions}%
{Per què, en dividir o multiplicar per un nombre negatiu, la desigualtat gira de sentit.}

\subsection{Idea clau}

Resoldre una inequació s'assembla molt a resoldre una equació: aïllem la $x$. Però
hi ha \textbf{un parany}: quan multipliquem o dividim els dos costats per un nombre
\textbf{negatiu}, el sentit de la desigualtat \textbf{gira}. Aquesta és la regla que
més errors causa, i avui la descobrirem amb un experiment.

\begin{idea}
\textbf{Les tres regles per transformar una inequació} (sense canviar-ne la solució):
\begin{enumerate}[leftmargin=1.6em, itemsep=2pt]
  \item \textbf{Sumar o restar} el mateix a tots dos costats: el sentit
        \textbf{es manté}.
  \item \textbf{Multiplicar o dividir} per un nombre \textbf{positiu}: el sentit
        \textbf{es manté}.
  \item \textbf{Multiplicar o dividir} per un nombre \textbf{negatiu}: el sentit
        \textbf{gira} ($<$ passa a $>$, $\le$ passa a $\ge$, i al revés).
\end{enumerate}
\textbf{Per què gira?} Sabem que $2<5$. Si multipliquem per $-1$: $-2$ i $-5$. Però
$-2>-5$! La desigualtat s'ha girat: per això, en operar amb negatius, cal capgirar
el símbol.
\end{idea}

\subsection{Exemples}

\begin{exemple}[l'experiment del conflicte]
Resolem $-2x+10\le 4$. Aïllem el terme amb $x$ restant $10$:
\[
-2x\le 4-10 \;\Rightarrow\; -2x\le -6.
\]
Ara dividim per $-2$. Com que és \textbf{negatiu}, \textbf{girem} el símbol:
\[
x\ge \frac{-6}{-2} \;\Rightarrow\; x\ge 3.
\]
\textbf{Comprovació:} provem $x=5$ (que compleix $x\ge 3$) a l'original:
$-2\cdot 5+10=0\le 4$, cert. Si \emph{no} haguéssim girat el símbol, hauríem escrit
$x\le 3$, i $x=0$ donaria $10\le 4$, fals. Per tant, girar és imprescindible.
\end{exemple}

\begin{exemple}[amb parèntesis]
Resolem $3(x-2)-5x\le 4x+6$. Primer desfem el parèntesi i agrupem:
\[
3x-6-5x\le 4x+6 \;\Rightarrow\; -2x-6\le 4x+6.
\]
Passem les $x$ a un costat i els números a l'altre:
\[
-2x-4x\le 6+6 \;\Rightarrow\; -6x\le 12.
\]
Dividim per $-6$ (negatiu) i \textbf{girem}:
\[
x\ge \frac{12}{-6} \;\Rightarrow\; x\ge -2.
\]
\end{exemple}

\subsection{Exercicis resolts}

\begin{exresolt}[sense canvi de signe]
Resol $2x+3>11$ i escriu la solució en forma d'interval.
\solucio
\[
2x>11-3 \;\Rightarrow\; 2x>8 \;\Rightarrow\; x>\frac{8}{2} \;\Rightarrow\; x>4.
\]
Dividim per $2$ (positiu): el símbol \emph{no} gira. Solució: $x>4$, interval
$(4,+\infty)$.
\resposta{$x>4$, és a dir $(4,+\infty)$.}
\end{exresolt}

\begin{exresolt}[amb canvi de signe i comprovació]
Resol $5-3x\le 2$, escriu la solució en interval i comprova-la.
\solucio
\[
-3x\le 2-5 \;\Rightarrow\; -3x\le -3.
\]
Dividim per $-3$ (negatiu) i \textbf{girem}:
\[
x\ge \frac{-3}{-3} \;\Rightarrow\; x\ge 1.
\]
Interval $[1,+\infty)$. \textbf{Comprovació} amb $x=2$ (compleix $x\ge 1$):
$5-3\cdot 2=-1\le 2$, cert.
\resposta{$x\ge 1$, interval $[1,+\infty)$.}
\end{exresolt}

\subsection{Exercicis per fer}

\noindent\textit{Per a cadascuna: resol, escriu la solució en forma d'interval i
comprova-la amb un valor.}

\begin{exercicis}
\begin{exlist}
  \item $x+5>9$
  \item $3x-2\le 10$
  \item $-x+4\ge 1$ \quad (compte amb el signe!)
  \item $-2x-1<5$ \quad (compte amb el signe!)
  \item $4x+3\ge 2x+11$
  \item \textbf{(Amb parèntesis)} $2(x-3)\le x-1$
  \item \textbf{(Caça l'error)} Algú resol $-4x>8$ així: «$x>\dfrac{8}{-4}=-2$, per
        tant $x>-2$». On és l'error? Resol-la correctament.
\end{exlist}
\end{exercicis}

% ---------------------------------------------------------------------
\fullsolucions{Sessió 6}
\begin{solucions}
\begin{sollist}
  \item $x>4$, interval $(4,+\infty)$.
  \item $3x\le 12 \Rightarrow x\le 4$, interval $(-\infty,4]$.
  \item $-x\ge -3$; dividim per $-1$ i girem: $x\le 3$, interval $(-\infty,3]$.
  \item $-2x<6$; dividim per $-2$ i girem: $x>-3$, interval $(-3,+\infty)$.
  \item $4x-2x\ge 11-3 \Rightarrow 2x\ge 8 \Rightarrow x\ge 4$, interval $[4,+\infty)$.
  \item $2x-6\le x-1 \Rightarrow 2x-x\le -1+6 \Rightarrow x\le 5$, interval $(-\infty,5]$.
  \item L'error és \textbf{no girar} el símbol en dividir per $-4$ (negatiu). El
        correcte: $-4x>8 \Rightarrow x<\dfrac{8}{-4} \Rightarrow x<-2$, interval
        $(-\infty,-2)$.
\end{sollist}
\end{solucions}
